\subsection{An exact query characterization with exponential size}
\label{sec:query-flow-lp}

For completeness, an instance has an exact finite optimization
description if exponential preprocessing is allowed.

Let \(p:\{-1,1\}^m\to[0,1]\) be the required probability of
returning one.  Write \(\Sigma=\{-1,*,1\}^m\) for the partial
assignments, and \(\sigma\preceq x\) if \(x\) extends
\(\sigma\).  At a partial assignment \(\sigma\), introduce
nonnegative variables \(u_\sigma,v_\sigma\) for stopping with
one or zero and \(a_{\sigma i}\) for querying an unassigned
coordinate \(i\).  Define the incoming flow by
\[
 w_{*^m}=1,\qquad
 w_\sigma=
 \sum_{i:\sigma_i\ne *}a_{\sigma^{(i\leftarrow *)},i}
 \quad(\sigma\ne *^m).
\]
Minimize \(C\), subject to
\begin{align*}
 w_\sigma&=u_\sigma+v_\sigma+
                 \sum_{i:\sigma_i=*}a_{\sigma i}
 &&(\sigma\in\Sigma),\\
 \sum_{\sigma\preceq x}u_\sigma&=p(x)
 &&(x\in\{-1,1\}^m),\\
 \sum_{\sigma\preceq x}\sum_{i:\sigma_i=*}a_{\sigma i}
 &\le C
 &&(x\in\{-1,1\}^m).
\end{align*}
Denote the value of this program by \(\mathsf{Flow}(p)\).
It has \(O(m3^m)\) variables and \(O(3^m+2^m)\) constraints.

\begin{lemma}[Flow characterization]\label{lem:instanceflow}
The infimum of the worst-input expected number of distinct probes
among exact samplers for \(p\) is \(\mathsf{Flow}(p)\).
This statement counts queries only.  It permits exponential
preprocessing and storage.
\end{lemma}

\paragraph{Proof idea.}
Merge decision-tree histories that reveal the same partial input.
Their stopping and query probabilities form a conserved flow. Conversely,
normalizing a feasible flow gives a sampler; a small complete-scan correction
then shows that computable probabilities attain the same infimum.

\begin{proof}
First consider a randomized decision tree with no repeated query.
All histories that end at the same partial assignment \(\sigma\)
have probabilities determined by the answers recorded in \(\sigma\).
Their probabilities are therefore the same on every full input
extending \(\sigma\).  Add their probabilities of stopping with
one, stopping with zero, or querying \(i\), to obtain
\(u_\sigma,v_\sigma,a_{\sigma i}\).
The incoming-flow identity and conservation hold by construction.
On a full input \(x\), summing the one-stopping flows gives the
output probability, and summing the query flows gives the expected
query count.  Every exact sampler supplies a feasible solution with
\(C\) equal to its worst-input expected query count.  The same
finite-input reduction used for the transcript lemma applies to
almost surely terminating algorithms with arbitrary internal work.

Conversely, at a state with \(w_\sigma>0\), stop with one, stop
with zero, or query \(i\), with respective probabilities
\(u_\sigma/w_\sigma\), \(v_\sigma/w_\sigma\), and
\(a_{\sigma i}/w_\sigma\).  A state with zero incoming flow is
unreachable.  Each query fixes one new coordinate, so the process
stops after at most \(m\) queries.  Induction on the number of
fixed coordinates shows that its probability of reaching
\(\sigma\), on any input extending \(\sigma\), is exactly
\(w_\sigma\).  The output and cost constraints thus give the
specified law and cost.  This proves equality when arbitrary
real transition probabilities are permitted.

For the feed-forward networks in this paper, the same infimum is obtained by
computable finite-bit samplers.  Their output probabilities belong
to \([\varepsilon,1-\varepsilon]\) for an explicitly computable
\(\varepsilon>0\).  For example, any rational number below
\((1-\tanh(1+s))/2\) suffices when the output row norm is at
most \(s\).  All \(2^m\) probabilities are computable to
arbitrary precision by evaluating the network.

The flow value is continuous quantitatively on this interior cube.
If \(p,q\in[\varepsilon,1-\varepsilon]^{2^m}\) and
\(\delta=\|p-q\|_\infty\le\varepsilon\), put
\(\gamma=\delta/\varepsilon\).  For \(\delta>0\), the table
\[
 r(x)=\frac{p(x)-(1-\gamma)q(x)}{\gamma}
\]
belongs to \([0,1]^{2^m}\).  Mix a sampler for \(q\), with
probability \(1-\gamma\), and a complete input scan followed by
an \(r(x)\)-coin, with probability \(\gamma\).  It follows that
\[
 |\mathsf{Flow}(p)-\mathsf{Flow}(q)|
 \le\frac{m}{\varepsilon}\|p-q\|_\infty.
\]
The case \(\delta=0\) is immediate.  A rational approximation
\(q\) admits an exactly solved rational flow program.  Using the
same mixture with a certified approximation error and
\(\varepsilon/2\) in place of \(\varepsilon\) yields an exact
finite-bit sampler whose query cost approaches
\(\mathsf{Flow}(p)\).  The residual scalar probability is
computable; a lazy fair-bit comparison samples it exactly.  Hence
computability does not change the infimum.
\end{proof}

The flow program accounts for all signed cancellations because it
uses the complete output table.  Its size exceeds the permitted
preprocessing budget exponentially.
